\section{Nonlinear parity, rigidity, and quantitative residual gaps}
\label{app:nonlinear-parity}

\subsection{A pointwise identity for every fixed-shape dilation}

First consider a centered scaling path $p_t(x)=g(t)^{-d}\rho(x/g(t))$,
and abbreviate $\alpha=g'/g$ and $\kappa=g''/g=\alpha'+\alpha^2$.
Every generating field can be written $u=\alpha x+r$ with
$\nabla\cdot(p_tr)=0$.  Distributionally,
\begin{align}
 \partial_t p_t&=-\alpha\nabla\cdot(p_tx),\nonumber\\
 \partial_{tt}p_t
 &=\alpha^2\nabla\cdot\nabla\cdot(p_txx^\top)
       -\kappa\nabla\cdot(p_tx).
 \label{eq:dilation-second-density}
\end{align}
For clarity, the double divergence of a tensor $S$ is
$\sum_{ij}\partial_i\partial_jS_{ij}$.  The identity
\[
 \nabla\cdot\nabla\cdot(p_txx^\top)
 =\nabla\cdot(p_tx)+\nabla\cdot[x\nabla\cdot(p_tx)]
\]
proves the second formula directly.  Differentiating the continuity
equation of $u$ gives
\[
 \partial_{tt}p_t
 =\nabla\cdot\nabla\cdot(p_tuu^\top)
      -\nabla\cdot(p_t\acc_u).
\]
Moreover
$\nabla\cdot\nabla\cdot(p_txr^\top)
=\nabla\cdot(p_tr)=0$, and the transposed cross tensor has the same
double divergence.  Acceleration preservation therefore implies
\begin{equation}
 \nabla\cdot\nabla\cdot(p_trr^\top)
 =-\kappa\nabla\cdot(p_tx).
 \label{eq:dilation-stress}
\end{equation}
All identities follow either pointwise with the stated smoothness or
by testing against compactly supported smooth functions.

Fix time and write $p\propto e^{-V_t}$.  Since
$\operatorname{div}r=\nabla V_t\cdot r$,
\begin{align*}
 p^{-1}\operatorname{div}\operatorname{div}(prr^\top)
 &=p^{-1}\operatorname{div}(p(Dr)r)\\
 &=\operatorname{tr}((Dr)^2)
   +r\cdot\nabla(\operatorname{div}r)
   -\nabla V_t\cdot(Dr)r\\
 &=\operatorname{tr}((Dr)^2)+r^\top D^2V_t\,r.
\end{align*}
Dividing \eqref{eq:dilation-stress} by $p$ proves
\eqref{eq:nonlinear-parity-identity}.  The trace is that of the matrix
square, not the squared Frobenius norm; it can be negative.

\subsection{Large level sets and the all-generator rigidity theorem}

\begin{proof}[Proof of Theorem~\ref{thm:nonlinear-odd-rigidity}]
The tail condition implies Gaussian upper tails.  In fact, on every ray
outside $R_0$, $\partial_rV(re)\ge cr$; integration gives
$V(re)\ge cr^2/2-C$ uniformly in unit $e$.  Thus all polynomial moments
are finite.  The affine-mean restriction already proved in
Corollary~\ref{cor:mean-action} gives $m''=0$.  Removing this affine mean
does not change material acceleration or its weighted divergence, so
we work with $m=0$.

Let $M_\rho=\E_\rho\|Z\|^2>0$.  The gradient base field $v=\alpha x$
belongs to the weighted-gradient closure, and $r=u-v$ is orthogonal to
it.  The second-moment identity therefore gives
\begin{equation}
 \|r_t\|_{L^2(p_t)}^2
 =\tfrac12\frac{\dd^2}{\dd t^2}(g(t)^2M_\rho)
      -(g'(t))^2M_\rho
 =g(t)g''(t)M_\rho.
 \label{eq:scale-gauge-energy}
\end{equation}
Hence $g''\ge0$, and $g''=0$ forces $r=0$ almost everywhere, then
everywhere by continuity and positivity of $p_t$.

Suppose $g''>0$ at some time.  Freeze that time, let
$L=\operatorname{Lip}(r)$, and put $W(x)=V(x/g)$, dropping an irrelevant
additive constant.  For sufficiently large $\|x\|$ there are positive
$c_t,M_t$ with
$x\cdot\nabla W\ge c_t\|x\|^2$ and $D^2W\preceq M_tI$.
Choose $R_1$ beyond this region and a level
$s>\max_{\|x\|\le R_1}W(x)$.
On every ray from zero, $W(re)$ is strictly increasing for $r\ge R_1$
and tends to infinity.  It therefore crosses $s$ at exactly one radius
$R_s(e)>R_1$.  The implicit function theorem shows that $R_s$ is smooth.
Thus $\{W=s\}$ is a compact smooth radial graph diffeomorphic to
$S^{d-1}$.  As $s\to\infty$, its minimum Euclidean radius tends to
infinity.  Also $\|\nabla W(x)\|\ge c_t\|x\|$ there.

Project $r$ onto the tangent bundle of this level set.  Because $d-1$
is even, the hairy-ball theorem~\citep[Section~2.2]{hatcher2002algebraic}
gives a point $x_s$ where that projection vanishes.  For $d=1$, the
tangent space of a level point is zero and the same conclusion is
immediate.  At $x_s$, $r$ is parallel to $\nabla W$.  Weighted divergence
then gives
\begin{equation}
 \|r(x_s)\|
 =\frac{|\nabla W(x_s)\cdot r(x_s)|}{\|\nabla W(x_s)\|}
 =\frac{|\operatorname{div}r(x_s)|}{\|\nabla W(x_s)\|}
 \le\frac{dL}{\|\nabla W(x_s)\|}.
 \label{eq:levelset-pole}
\end{equation}
But the pointwise identity requires
\[
 \kappa(x_s\cdot\nabla W(x_s)-d)
 \le dL^2+
      \frac{M_td^2L^2}{\|\nabla W(x_s)\|^2}.
\]
The left side tends to infinity and the right side remains bounded.
This contradiction proves that $g''=0$ everywhere.  Equation
\eqref{eq:scale-gauge-energy} proves uniqueness of the generator.

Conversely, if $m$ and positive $g$ are affine, the field
\eqref{eq:odd-unique-generator} has zero material acceleration.  Its
exact flow from $s$ to $t$ is
$m(t)+(g(t)/g(s))(x-m(s))$; since both scalar and mean increments are
affine, this map equals its Euler step on every interval.  Its global
regularity and moment assumptions follow from $\min_{[0,T]}g>0$ and
the source moments.
\end{proof}

The topological fact used here can be seen directly.  A nowhere-zero
continuous tangent field on $S^{d-1}$ would normalize to $\tau(e)$ with
$\|\tau(e)\|=1$ and $e\cdot\tau(e)=0$.  Then
$H_s(e)=\cos(\pi s)e+\sin(\pi s)\tau(e)$ is a sphere-valued homotopy
from the identity to the antipodal map.  Their degrees are $1$ and
$(-1)^d$, which differ for odd $d$.  This argument is classical; the
new use is to force an obstruction on the unconstrained nonlinear
generator through its weighted-divergence and acceleration identities.

\paragraph{Finite Gaussian mixtures.}
Let $\rho=\sum_j\pi_j\rho_j$ with nondegenerate Gaussian components,
and set $w_j=\pi_j\rho_j/\rho$, $Q_j=\Sigma_j^{-1}$,
$s_j(x)=Q_j(x-\mu_j)$.  Differentiating the logarithm yields
$\nabla V=\sum_jw_js_j$ and
$D^2V=\sum_jw_jQ_j-\operatorname{Cov}_w(s_j)$.
The covariance matrix is positive semidefinite, so
$D^2V\preceq\Lambda I$.  The displayed radial lower bound in the
main text gives $x\cdot\nabla V\ge(\lambda/2)\|x\|^2$ whenever
$\|x\|\ge2B/\lambda$.  These prove every hypothesis of
Theorem~\ref{thm:nonlinear-odd-rigidity}.  No separation, equal covariance,
radial symmetry, or log-concavity of the mixture is required.

\paragraph{Gaussian specialization.}
For $p=\mathcal N(0,g^2I)$, the pointwise identity reduces to
\begin{equation}
 \|r(x)\|^2+g^2\operatorname{tr}((Dr(x))^2)
 =\frac{g''}{g}\|x\|^2-dgg''.
 \label{eq:gaussian-pointwise-parity}
\end{equation}
When $g''>0$, bounded $Dr$ would force asymptotically nonzero tangential
speed on every sufficiently large sphere.  Equation
\eqref{eq:levelset-pole} is the topological contradiction.  This is not
the eigenvalue pairing argument for linear skew matrices.

\begin{proposition}[Necessary Jacobian growth and polynomial rigidity]
\label{prop:gaussian-escape-rigidity}
Fix $p=\mathcal N(0,g^2I_d)$, $g>0$, and $\kappa>0$.  Let $r$ be a
$C^2$ field satisfying $\operatorname{div}(pr)=0$ and the stress equation
\eqref{eq:dilation-stress}.  No global Lipschitz or polynomial growth
assumption is imposed here.  If $d$ is odd, then, for every $R>g\sqrt d$,
\begin{equation}
 \sup_{\|x\|=R}\|Dr(x)\|_{\rm op}
 \ge \frac{\sqrt\kappa\,R}{g\sqrt d}
       \sqrt{\frac{R^2-dg^2}{R^2+dg^2}}.
 \label{eq:gaussian-jacobian-growth}
\end{equation}
Consequently the liminf of this supremum divided by $R$ is at least
$\sqrt\kappa/(g\sqrt d)$.  In every dimension, if $r$ is polynomial,
then $r(x)=Kx$, where $K^\top=-K$ and $K^2=-\kappa I_d$.
Thus no polynomial correction is possible in odd dimension.
\end{proposition}

\begin{proof}
Write $M_R=\sup_{\|x\|=R}\|Dr(x)\|_{\rm op}$.  At a hairy-ball pole
$Re$ of this sphere, $r(Re)$ is parallel to $e$.  Weighted divergence
therefore gives $\|r(Re)\|\le dg^2M_R/R$.  The Gaussian identity
\eqref{eq:gaussian-pointwise-parity} implies
\[
 \kappa(R^2-dg^2)
 \le\frac{d^2g^4M_R^2}{R^2}+dg^2M_R^2.
\]
Solving for $M_R$ proves \eqref{eq:gaussian-jacobian-growth}.

If $r$ is a nonzero polynomial of degree $m\ge2$, the leading term of
$\|r\|^2$ has degree $2m$ and is a nonzero sum of squares.  The
derivative term in \eqref{eq:gaussian-pointwise-parity} has degree at
most $2m-2$, while the right side has degree two.  This is impossible,
so $r(x)=Kx+b$ is affine.  The polynomial identity
$\operatorname{tr}K=g^{-2}x\cdot(Kx+b)$ forces $K^\top=-K$ and $b=0$.
The quadratic terms in \eqref{eq:gaussian-pointwise-parity} then give
$K^\top K=\kappa I_d$, equivalently $K^2=-\kappa I_d$.  An invertible
real skew-symmetric matrix cannot have odd size.  These conclusions
are instantaneous; any time-dependent polynomial generator of the
Gaussian scaling path must meet them at every time.
\end{proof}

\subsection{Why all convex-observable tests can pass}

\begin{proof}[Proof of Corollary~\ref{cor:convex-hierarchy-insufficient}]
Let $X_t=g(t)Z$, with $Z$ centrally symmetric.  Differentiating along
the base flow $v=\alpha x$ gives
\[
 F_\phi''=\E[v^\top D^2\phi(X_t)v]
             +\kappa\E[X_t\cdot\nabla\phi(X_t)].
\]
Central symmetry shows
\[
 \E[X\cdot\nabla\phi(X)]
 =\tfrac12\E\bigl[X\cdot
       (\nabla\phi(X)-\nabla\phi(-X))\bigr]\ge0,
\]
because the gradient of a convex differentiable function is monotone.
For $g''\ge0$ the second term is nonnegative.  The first term is at
least $K_\phi$, since the base field is one admissible velocity in that
infimum.  This proves all the convex and weighted-energy inequalities,
under the stated integrability for each test.

If $C_\rho=\operatorname{Cov}(Z)$ is positive definite, then
$C=g^2C_\rho$ and the choice $B=gg'C_\rho$ in the covariance test has
Schur complement $gg''C_\rho\succeq0$.  For singular $C_\rho$, the
same positive semidefinite block matrix follows by a direct tensor
factorization on its support.  The odd-dimensional exponential Gaussian
has strictly convex scale, so Theorem~\ref{thm:nonlinear-odd-rigidity}
nevertheless excludes every admissible generator.
\end{proof}

\subsection{Conditional variance identities and the Gaussian residual floor}

Fix a time for a Gaussian scaling path and write $g=g(t)$,
$\alpha=g'/g$, $\kappa=g''/g$, $c_t=gg''$.  Let $X\sim p_t$ and
$r=u-\alpha x$.  For every unit vector $e$, write $Y=e\cdot X$,
$q_g(y)=(2\pi g^2)^{-1/2}e^{-y^2/(2g^2)}$, and
$Q_e(y)=\E[(e\cdot r(X))^2\mid Y=y]$.
Projecting $\operatorname{div}(p_tr)=0$ onto this coordinate gives
\begin{equation}
 \E[e\cdot r(X)\mid Y=y]=0\quad\text{for every }y.
 \label{eq:projected-gauge-mean}
\end{equation}
Indeed the derivative of $q_g(y)$ times this conditional mean is zero;
the only constant compatible with Gaussian tails and linear field growth
is zero.  The conditional mean and variance have continuous versions,
by dominated convergence and Lipschitz growth, so these conclusions are
available at every chosen $y$, not merely almost everywhere.

For any twice differentiable scalar test with bounded first and second
derivatives, the scaling identity and the physical chain rule give
\begin{equation}
 \int\nabla[\phi(e\cdot x)]\cdot\acc_u(x)\,p_t(\dd x)
 =\int_{\R}\phi''(y)\{c_t-Q_e(y)\}q_g(y)\,\dd y.
 \label{eq:gaussian-directional-defect}
\end{equation}
To verify this, if $s(t)=g(t)^2$ then
$\partial_tq_g=(s'/2)\partial_{yy}q_g$ and
$\partial_{tt}q_g=(s''/2)\partial_{yy}q_g+
(s'^2/4)\partial_{yyyy}q_g$.
The latter equals
$\partial_{yy}[q_g(\alpha^2y^2+c_t)]$.
By \eqref{eq:projected-gauge-mean}, the conditional squared velocity is
$\alpha^2y^2+Q_e(y)$.  Subtracting the velocity-Hessian term from
the second-derivative chain rule proves
\eqref{eq:gaussian-directional-defect}.  If acceleration preservation
were imposed, this also forces $Q_e(y)=c_t$: a function whose second
distributional derivative is zero is affine, and integrability of
$q_g(Q_e-c_t)$ excludes nonzero affine terms.

\begin{proof}[Proof of Theorem~\ref{thm:quantitative-odd-floor}]
Let $L=\operatorname{Lip}(r)$.  Pointwise weighted divergence gives
\begin{equation}
 |x\cdot r(x)|=g^2|\operatorname{div}r(x)|\le C,
 \qquad C=dg^2L.
 \label{eq:gaussian-radial-flux-bound}
\end{equation}
For every radius $R>0$, the tangential projection of $r$ on the sphere
of radius $R$ has a zero $Re_R$.  Hence $r(Re_R)$ is parallel to $e_R$.
For $y=R+s$, $0\le s\le g$, condition on $e_R\cdot X=y$ and write
$X=ye_R+Z_\perp$, where
$Z_\perp\sim\mathcal N(0,g^2I_{d-1})$.  Since
$Z_\perp\cdot r(Re_R)=0$, \eqref{eq:gaussian-radial-flux-bound} and
Lipschitzness give
\[
 |e_R\cdot r(ye_R+Z_\perp)|
 \le\frac{C+L\|Z_\perp\|\sqrt{s^2+\|Z_\perp\|^2}}{R}.
\]
Minkowski's inequality, $\E\|Z_\perp\|^2=(d-1)g^2$, and
$\E\|Z_\perp\|^4=(d-1)(d+1)g^4$ imply
\begin{equation}
 \sup_{R\le y\le R+g}\sqrt{Q_{e_R}(y)}
 \le\frac{Lg^2}{R}\{d+\sqrt{d^2+d-2}\}
 =\frac{Lg^2A_d}{R}.
 \label{eq:gaussian-pole-strip}
\end{equation}
Choose $R=g\max\{1,2A_dL/\sqrt\kappa\}$.  Then
$Q_{e_R}\le c_t/4$ throughout this strip.

Choose a scalar convex $C^{1,1}$ test with
$\phi''=g^{-1}\mathbf1_{[R,R+g]}$, and $\phi'=0$ below $R$, so
$0\le\phi'\le1$.  Formula~\eqref{eq:gaussian-directional-defect}
and the dual definition of the projected norm yield
\[
 \varepsilon_t\ge\frac{3c_t}{4g}\Pr\{R\le Y\le R+g\}
 \ge\frac{3g''}{4\sqrt{2\pi}}
       \exp[-(R/g+1)^2/2].
\]
The test is legitimate by smooth mollification and then compact cutoff:
its first and second derivatives are bounded, it grows at most linearly,
and its gradients converge in $L^2(p_t)$.  Gaussian moments dominate
the chain-rule terms.  Weak derivatives at the strip endpoints do not
create an additional term, since $\phi'$ is continuous.
Weakening $3/4$ to $1/2$ gives
\eqref{eq:quantitative-odd-floor}.  Solving this inequality for $L$
under the stated threshold on $\eta$ gives
\eqref{eq:odd-regularity-cost}.
\end{proof}

\subsection{A nonexplicit uniform gap for the broad source class}

\begin{proposition}[Positive residual gap at fixed Lipschitz bound]
\label{prop:compact-odd-gap}
Fix a positive density $p\propto e^{-W}$ satisfying
\eqref{eq:parity-tail-assumptions}, an odd $d$, $\kappa>0$, and a finite
$L\ge0$.  Among all globally $L$-Lipschitz fields $r$ with
$\operatorname{div}(pr)=0$, the quantity
\[
 \mathcal R_\kappa(r)=
 \sup_{\substack{\phi\in C_c^\infty\\\|\nabla\phi\|_{L^2(p)}\le1}}
 \left|\kappa\int x\cdot\nabla\phi\,\dd p
       -\int r^\top D^2\phi\,r\,\dd p\right|
\]
attains a finite strictly positive minimum.  For a generator of the
scaling path at this time, $\mathcal R_\kappa(r)=\|\Pi_p\acc_u\|$.
\end{proposition}

\begin{proof}
Cutoff tests in the weighted-divergence equation give $\E_pr=0$.
Consequently $\|r(0)\|\le L\E_p\|X\|$, giving a uniform linear-growth
bound for this class.  Arzela--Ascoli and a diagonal extraction give a
locally uniform subsequential limit of every sequence.  The Gaussian
tail bound and uniform linear growth upgrade this to strong $L^2(p)$
convergence.  Both the Lipschitz bound and weighted divergence survive
the limit, so the class is compact in $L^2(p)$.
For each fixed test $\phi$, its displayed quadratic integral is continuous
in this topology.  A supremum of continuous functions is lower
semicontinuous.  The minimum is attained, and is finite because the
candidate $r=0$ has value $\kappa\|x\|_{L^2(p)}$.

If the minimum were zero, its Lipschitz minimizer would satisfy
\eqref{eq:dilation-stress} distributionally.  The pointwise identity
extends to this regularity almost everywhere.  Here is a justification:
$r\in W^{1,\infty}_{\rm loc}$ and
$\operatorname{div}r=\nabla W\cdot r\in W^{1,\infty}_{\rm loc}$.
Mollification, the product rule, and convergence in local Sobolev norms
give
\[
 \operatorname{div}((Dr)r)
 =\operatorname{tr}((Dr)^2)+r\cdot\nabla(\operatorname{div}r)
\]
as distributions.  Adding the weight cancels the first-order score
terms exactly as above.  Since $\|Dr\|\le L$ almost everywhere and
$D^2W\preceq MI$ in the tails, this yields
\[
 \kappa(x\cdot\nabla W-d)\le dL^2+M\|r(x)\|^2
\]
almost everywhere there.  Both sides are continuous, so the inequality
holds everywhere.  Similarly
$|\nabla W\cdot r|\le dL$ holds everywhere.  Applying the level-set
hairy-ball argument to the continuous field $r$ gives the same
contradiction as in Theorem~\ref{thm:nonlinear-odd-rigidity}.

For an actual scaling generator, integrating its second-derivative
identity against $\phi$ gives the formula for $\mathcal R_\kappa$.
The cross term is zero because
$\int r\cdot D^2\phi\,x\,\dd p
=\int r\cdot\nabla(x\cdot\nabla\phi-\phi)\,\dd p=0$.
The dual projection formula then proves the asserted equality.
\end{proof}

\subsection{A separate radial-score extension}

For completeness, odd-dimensional affine-scale rigidity also holds for
radial densities that need not meet the upper-Hessian assumption.
Suppose $\rho>0$ is smooth and
\[
 \nabla\log\rho(z)=-b(\|z\|)z,\qquad
 0<\lambda\le b(r)\le\Lambda<\infty.
\]
These hypotheses give Gaussian tails.  At a fixed scale $g$, let $q_g$
be any one-dimensional marginal and define its scalar Stein kernel
\[
 \tau_g(y)=q_g(y)^{-1}\int_y^\infty s q_g(s)\,\dd s.
\]
The marginal is even and
$-q_g'(y)/q_g(y)=y\E[b(\|X\|/g)\mid e\cdot X=y]/g^2$.
Integration of these score bounds gives
$g^2/\Lambda\le\tau_g(y)\le g^2/\lambda$.

For a scale path, the projected continuity equation first gives
$\E[e\cdot r\mid e\cdot X=y]=0$.  Equation
\eqref{eq:dilation-stress} then implies
\[
 \E[(e\cdot r)^2\mid e\cdot X=y]=\kappa\tau_g(y).
\]
Indeed, writing the conditional variance as $Q_e$, the equation is
$(q_gQ_e)''=-\kappa(yq_g)'$; integrability removes the affine integration
constants, and $(q_g\tau_g)'=-yq_g$.
On the other hand, weighted divergence gives
$|x\cdot r(x)|\le dg^2L/\lambda$.
At a hairy-ball pole $Re_R$, tangency vanishes.  For
$X=Re_R+Z_\perp$ conditioned on $e_R\cdot X=R$, the same cancellation
$Z_\perp\cdot r(Re_R)=0$ bounds
\[
 |e_R\cdot r(X)|\le
 \frac{dg^2L/\lambda+L\|Z_\perp\|^2}{R}.
\]
Conditional integration by parts gives, uniformly in $R$,
$\E\|Z_\perp\|^2\le(d-1)g^2/\lambda$ and
$\E\|Z_\perp\|^4\le(d-1)(d+1)g^4/\lambda^2$.
For the latter identities, test the transverse score equation with
$z$ and $\|z\|^2z$ and use its coefficient lower bound
$\lambda/g^2$; cutoff limits are justified by the uniform Gaussian
tails.  Thus the conditional variance tends to zero as $R\to\infty$,
contradicting $\kappa\tau_g(R)\ge gg''/\Lambda>0$ if $g''>0$.
Equation~\eqref{eq:scale-gauge-energy} handles the other signs and
completes the same rigidity conclusion.  This argument uses no bound
on derivatives of $b$.
